\documentclass[tikz,border=6pt]{standalone}
\usepackage{ctex}
\usetikzlibrary{arrows.meta,calc}
% Shared visual language for LFS architecture diagrams.
\RequirePackage{../../mymath}

\definecolor{LFSBlue}{HTML}{2563EB}
\definecolor{LFSBlueSoft}{HTML}{DBEAFE}
\definecolor{LFSGreen}{HTML}{15803D}
\definecolor{LFSGreenSoft}{HTML}{DCFCE7}
\definecolor{LFSOrange}{HTML}{C2410C}
\definecolor{LFSOrangeSoft}{HTML}{FFEDD5}
\definecolor{LFSRed}{HTML}{B91C1C}
\definecolor{LFSRedSoft}{HTML}{FEE2E2}
\definecolor{LFSPurple}{HTML}{7E22CE}
\definecolor{LFSPurpleSoft}{HTML}{F3E8FF}
\definecolor{LFSGray}{HTML}{475569}
\definecolor{LFSGraySoft}{HTML}{F1F5F9}

\tikzset{
  lfs picture/.style={
    font=\small,
    node distance=8mm and 12mm,
    >=Latex,
    every path/.style={line cap=round, line join=round}
  },
  block/.style={
    draw=LFSBlue,
    fill=LFSBlueSoft,
    rounded corners=2pt,
    line width=0.8pt,
    align=center,
    inner xsep=7pt,
    inner ysep=5pt,
    minimum height=10mm
  },
  compute/.style={
    block,
    draw=LFSPurple,
    fill=LFSPurpleSoft
  },
  storage/.style={
    block,
    draw=LFSGreen,
    fill=LFSGreenSoft
  },
  interface/.style={
    block,
    draw=LFSOrange,
    fill=LFSOrangeSoft
  },
  risk/.style={
    block,
    draw=LFSRed,
    fill=LFSRedSoft
  },
  neutral/.style={
    block,
    draw=LFSGray,
    fill=LFSGraySoft
  },
  decision/.style={
    diamond,
    draw=LFSOrange,
    fill=LFSOrangeSoft,
    line width=0.8pt,
    align=center,
    inner sep=2pt,
    aspect=2
  },
  group/.style={
    draw=LFSGray,
    rounded corners=3pt,
    dashed,
    line width=0.7pt,
    inner sep=6pt
  },
  arrow/.style={
    -{Latex[length=2.4mm,width=1.6mm]},
    draw=LFSGray,
    line width=0.9pt
  },
  data arrow/.style={
    arrow,
    draw=LFSBlue
  },
  control arrow/.style={
    arrow,
    draw=LFSOrange,
    dashed
  },
  residual/.style={
    arrow,
    draw=LFSGreen,
    rounded corners=4pt
  },
  label text/.style={
    font=\scriptsize,
    text=LFSGray,
    fill=white,
    inner sep=1.5pt
  },
  title/.style={
    font=\bfseries,
    text=LFSGray,
    align=center
  }
}


\begin{document}
\begin{tikzpicture}[
  my picture,
  event/.style={text width=57mm,inner sep=0pt,font=\small},
  date/.style={fill=white,text=MyGray,font=\scriptsize,inner xsep=2pt,inner ysep=1pt},
  phase/.style={font=\small\bfseries,text=MyGray,fill=white,inner xsep=4pt}
]

\newcommand{\timelinedate}[2]{\node[date] at (0,#1) {#2};}
\newcommand{\methodnode}[3]{%
  \node[event,anchor=east,align=right] at (-0.95,#1) {%
    \textcolor{MyGreen}{\bfseries #2}\;{\color{MyGray}\scriptsize[#3]}%
  };%
  \draw[MyGreen,line width=0.7pt] (-0.86,#1)--(-0.72,#1);%
  \fill[MyGreen] (-0.72,#1) circle (1.25pt);%
}
\newcommand{\modelnode}[3]{%
  \node[event,anchor=west,align=left] at (0.95,#1) {%
    \textcolor{MyBlue}{\bfseries #2}\;{\color{MyGray}\scriptsize[#3]}%
  };%
  \draw[MyBlue,line width=0.7pt] (0.72,#1)--(0.86,#1);%
  \fill[MyBlue] (0.72,#1) circle (1.25pt);%
}
\newcommand{\recentmodelnode}[3]{%
  \node[event,anchor=west,align=left] at (0.95,#1) {%
    \textcolor{MyPurple}{\bfseries #2}\;{\color{MyGray}\scriptsize[#3]}%
  };%
  \draw[MyPurple,line width=0.7pt] (0.72,#1)--(0.86,#1);%
  \draw[MyPurple,line width=0.8pt,fill=white] (0.72,#1) circle (1.25pt);%
}
\newcommand{\lanehead}[1]{%
  \node[font=\small\bfseries,text=MyGreen,anchor=east] at (-0.95,#1) {方法与架构};%
  \node[font=\scriptsize,text=MyGray] at (0,#1) {时间};%
  \node[font=\small\bfseries,text=MyBlue,anchor=west] at (0.95,#1) {基础模型};%
}

% 第一阶段：架构与预训练
\draw[MyGray!35,line width=0.6pt] (-6.7,0.92)--(6.7,0.92);
\node[phase] at (0,0.92) {架构与预训练（2017--2022）};
\lanehead{0.48}
\draw[MyGray!55,line width=0.8pt] (0,0.20)--(0,-5.18);

\timelinedate{-0.10}{2017.06}
\methodnode{-0.10}{Transformer}{论文}

\timelinedate{-0.82}{2018.06}
\modelnode{-0.82}{GPT}{开放权重}

\timelinedate{-1.54}{2018.10}
\modelnode{-1.54}{BERT}{开放权重}

\timelinedate{-2.26}{2019.02}
\modelnode{-2.26}{GPT-2}{分阶段开放}

\timelinedate{-2.98}{2020.01}
\methodnode{-2.98}{神经语言模型缩放定律}{论文}

\timelinedate{-3.70}{2020.05}
\modelnode{-3.70}{GPT-3}{API}

\timelinedate{-4.42}{2022.01}
\methodnode{-4.42}{InstructGPT / RLHF}{论文}

\timelinedate{-5.14}{2022.03}
\methodnode{-5.14}{Chinchilla缩放规律}{论文}

% 第二阶段：多模态与推理模型
\draw[MyGray!35,line width=0.6pt] (-6.7,-5.93)--(6.7,-5.93);
\node[phase] at (0,-5.93) {多模态与推理模型（2023--2026）};
\lanehead{-6.37}
\draw[MyGray!55,line width=0.8pt] (0,-6.65)--(0,-15.40);

\timelinedate{-6.94}{2023.02}
\modelnode{-6.94}{LLaMA}{研究权重}

\timelinedate{-7.66}{2023.03}
\modelnode{-7.66}{GPT-4 · Claude 1}{API}

\timelinedate{-8.38}{2024.02}
\modelnode{-8.38}{Gemini 1.5}{API}

\timelinedate{-9.10}{2024.03}
\modelnode{-9.10}{Claude 3}{API}

\timelinedate{-9.82}{2024.05}
\modelnode{-9.82}{DeepSeek-V2 · GPT-4o}{权重 / API}

\timelinedate{-10.54}{2024.09}
\methodnode{-10.54}{推理时计算}{训练范式}
\modelnode{-10.54}{OpenAI o1}{API}

\timelinedate{-11.26}{2024.12}
\modelnode{-11.26}{DeepSeek-V3}{开放权重}

\timelinedate{-11.98}{2025.01}
\methodnode{-11.98}{可验证推理强化学习}{公开方法}
\modelnode{-11.98}{DeepSeek-R1}{开放权重}

\timelinedate{-12.70}{2025.02}
\modelnode{-12.70}{Claude 3.7 Sonnet}{API}

\timelinedate{-13.42}{2025.05}
\modelnode{-13.42}{Claude 4}{API}

\timelinedate{-14.14}{2025.08}
\modelnode{-14.14}{GPT-5}{API}

\timelinedate{-14.86}{2026.09.01}
\recentmodelnode{-14.86}{Claude Fable/Mythos 5.1}{API / 受控}

\timelinedate{-15.58}{2026.09.10}
\recentmodelnode{-15.58}{DeepSeek-V4.1-Flash}{开放权重}

\node[font=\scriptsize,text=MyGray,anchor=west] at (-6.65,-16.26) {%
  \textcolor{MyGreen}{\rule{6pt}{6pt}} 方法与架构\quad
  \textcolor{MyBlue}{\rule{6pt}{6pt}} 基础模型\quad
  \textcolor{MyPurple}{$\circ$} 新近节点\quad
  方括号标注发布形态
};

\end{tikzpicture}
\end{document}
